\section{希格斯分支与库仑分支}
\label{sec:Higgs-and-Coulomb}

本节聚焦于库仑分支与希格斯分支的一个格外重要的情形：映入 $\mathbb{H}^n$ 的 $d=3$、$\mathcal{N}=4$ 超对称 $\sigma$ 模型，对它按子群 $G_c\subset U(n,\mathbb{H})$ 的作用做规范；其中的场常称为超多重态（hypermultiplet）。对应于映入 $\mathbb{H}^n$ 的映射的那些场，常被称为该理论的「物质内容」（matter content）。

通常更方便的做法是忘掉 $\mathbb{H}^n$ 上的坐标，而把它看成一个一般的 $\mathbb{H}$-模 $X$，并选取一个范数以及 $G_c$ 的一个作用。把 $X$ 视为带复结构 $I$ 的 $\C$-向量空间、并带上 $G_c$ 的复化 $G$ 的诱导作用，也会使我们的讨论简化；四元数 $J$ 与 $K$ 的作用可以编码在全纯辛形式 $\Omega(x,y)=\langle Jx,y\rangle +i\langle Kx,y\rangle$ 之中。

我们把相应的理论记为 $\mathcal{T}(X,G)$，并把希格斯分支与库仑分支记为 $\mathfrak{M}_{\operatorname{A/B}}(X,G)$。

这两个簇都有具体的数学描述，我们将在有限的篇幅内尽量加以说明。二者都可以从无穷维几何中的操作推出，所依据的原则是：一个物理理论的希尔伯特空间，由该理论相空间的几何量子化得到。当 $X=T^*N$、$N$ 是 $G$-表示时，这种几何量子化最容易处理。极其粗略地说，这个相空间来自从 $S^2$ 到商 $N/G$ 的余切丛、满足某些性质的映射。要解释它们，最容易的办法是把我们的 $S^2$ 形变成柱体
\[\mathsf{Cyl}_\delta=\{(x,y,z) \mid x^2+y^2\leq \delta, |z|\leq \delta\}\]
\begin{wrapfigure}{R}{.14\textwidth}
\begin{tikzpicture}[scale=.5,MyPersp,font=\large]
  % 底圆是 Oxy 平面内的单位圆
  \def\h{1.5}% 椭圆中心的高度（在柱体的轴线上）
    \def\a{-50}% 控制坐标轴的角度
\draw[->] (0,0,{.5*\h}) -- ({2*cos(\a)},{2*sin(\a)},{.5*\h}) node[left] {$x$};
\draw[->] (0,0,{.5*\h}) -- ({2*sin(\a)},{-2*cos(\a)},{.5*\h}) node[right] {$y$};
\draw[->] (0,0,{.5*\h}) -- (0,0,{1.3*\h}) node[above] {$z$};
		% 背面填充
		\fill[blue,fill opacity=.05]  
		 (1,0,{\h})--(1,0,0)
		\foreach \t in {0,-2,-4,...,-180}
			{--({cos(\t)},{sin(\t)},0)}
-- (-1,0,0)--(-1,0,{\h})
		\foreach \t in {180,178,...,0}
			{--({cos(\t)},{sin(\t)},{\h})}--cycle;
		% 底圆（后半）
		\draw[gray,   very thick] (1,0,0)
		\foreach \t in {0,2,4,...,180}
			{--({cos(\t)},{sin(\t)},0)};
   % 柱体的正面	
		\fill[blue!20!white,fill opacity=.5]
		 (1,0,{\h})--(1,0,0)
		\foreach \t in {0,-2,-4,...,-180}
			{--({cos(\t)},{sin(\t)},0)}
-- (-1,0,0)--(-1,0,{\h})
		\foreach \t in {-180,-178,...,0}
			{--({cos(\t)},{sin(\t)},{\h})}--cycle;
   % 柱体的侧面
\draw[gray] (1,0,0)--(1,0,{\h});
\draw[gray] (-1,0,0)--(-1,0,{\h});
% 底圆（前半）
	\draw[gray, very thick] (1,0,0) 
		\foreach \t in {0,-2,-4,...,-180}
			{--({cos(\t)},{sin(\t)},0)};
			% 上圆
	\draw[gray, very thick] (1,0,\h) 
		\foreach \t in {2,4,...,360} 
			{--({cos(\t)},{sin(\t)},{\h})}--cycle;
                      \end{tikzpicture}
\end{wrapfigure}
的边界，其中 $\delta>0$ 为实数。

我们会经常提到这个柱体的顶面、底面与侧面，它们分别指 $z=\delta,-\delta$ 平面中的单位圆盘，以及夹在二者之间的那部分边界。
\begin{enumerate}[label=\Alph*., ref=\Alph*]
     \item \label{fns:A} 代数 $\alg_{\operatorname{A}}$ 是柱体到 $N/G$ 的映射空间上\emph{局部常值}函数构成的代数，这些映射在侧面上为常值，在顶面与底面上\emph{全纯}。
    \item \label{fns:B} 代数 $\alg_{\operatorname{B}}$ 是柱体到 $N/G$ 的映射空间上\emph{全纯}函数构成的代数，这些映射在侧面上为常值，在顶面与底面上\emph{局部常值}。
\end{enumerate}
我们这样表述，是为了突出其平行性，也就是说，A 扭转与 B 扭转的差别如何通过「局部常值」与「全纯」这两个词的位置体现出来。在下文中，我们会更仔细地展开对表述 (\ref{fns:A}) 与 (\ref{fns:B}) 中诸概念的解释，因为某种推广是必要的。

\subsection{希格斯分支}

首先考虑 B 扭转。尽管按字母顺序排在第二，相应的希格斯分支却更容易精确理解，因此在数学文献中一般也受到更多关注。按照表述 (\ref{fns:B})，代数 $\alg_{\operatorname{B}}$ 应当是常值映射 $S^2\to N/G$ 上的函数。这里，除了 $N/G$ 中的一个点，还应当考虑该商的一个余向量（见 \cite[\S 7.14]{bravermanCoulombBranches2018}）。

利用 $G$ 在辛向量空间 $(X,\Omega)$ 上的矩映射 $\mu\colon X\to \mathfrak{g}^*$，可以把上述代数描述得更具体。若 $X=T^*N$，则 $\mu^{-1}(0)$ 由所有这样的对构成：$n\in N$，以及一个在过 $n$ 的轨道的切空间上取零值的余向量 $\xi$（因而可以看作商上的余向量）。
\begin{definition}
	希格斯分支 $\mathfrak{M}_{\operatorname{B}}(X,G)$ 定义为一个{\bf 全纯辛商}，也就是说，有 $\alg_{\operatorname{B}}(X,G)=Z_{\operatorname{B}}(S^2)=\C[\mu^{-1}(0)]^G$，即 $\mu^{-1}(0)$ 上 $G$-不变的复多项式函数构成的代数，并且 $\mathfrak{M}_{\operatorname{B}}=\Spec\alg_{\operatorname{B}}$。该空间的点与 $\mu^{-1}(0)$ 中的闭 $G$-轨道一一对应。
\end{definition}

所得的簇通常带有辛奇点，因为我们施加的是某种版本的辛约化。人们自然会问：这个簇是否有辛消解？并非所有情形都如此，或许连大多数情形都不是，但在某些情形下确实有。

本文中我们将着重讨论两类特别的例子：阿贝尔规范理论与箭图规范理论。在这两种情形中，靶空间 $X=T^*N$ 都是 $G$ 的某个 $\C$-表示 $N$ 的余切丛。

\subsubsection{阿贝尔/超环面规范理论}


设 $G$ 是阿贝尔群。由于它连通且约化（即某个紧群的复化），这意味着对某个 $k$ 有 $G\cong (\C^{\times})^k$。对 $G$ 的任意 $\C$-表示 $N$，我们可以选取同构 $N\cong \C^n$，使得 $G\hookrightarrow D$ 是 $n\times n$ 对角矩阵全体构成的群 $D$ 的子群（也就是说，我们选取了一组权基）。

这些材料通常用于环面簇的构造：在矩映射的任一正则值处，辛向量空间 $(N,\omega_I)$ 的辛约化——等价地，GIT 商 $N/\!/_{\lambda}G$——将给出关于商 $F=D/G$ 的作用的拟射影环面簇。

于是，这个理论的希格斯分支 $\mathfrak{M}_{\operatorname{B}}(T^*N,G)$ 的构造，就是环面簇构造的四元数版本。所得的簇称为{\bf 超环面簇}或{\bf 环面超凯勒簇}。该簇的复维数为 $2(\dim N-\dim G)$。读者最熟悉的例子大概是下面这些：
\begin{yizhu}[超环面簇]
超环面簇是环面簇构造在超凯勒情形下的类比：把「辛向量空间取辛约化」换成「$T^*N$ 取超凯勒商」。它的组合数据来自 $G$ 在 $N$ 上的权所确定的超平面配置，所以组合问题往往能反过来读成超平面配置的组合问题。
\end{yizhu}
\begin{ex}\label{ex:TPn}
    设 $G=\{\varphi I \mid \varphi \in \C^{\times}\}\subset D$ 为 $n\times n$ 数量矩阵构成的子群。我们可以把 $T^*N$ 的元素看作一对向量：$n\times 1$ 列向量 $\mathbf{a}$ 与 $1\times n$ 行向量 $\mathbf{b}$；群 $G\cong \C^{\times}$ 通过
\[\varphi \cdot (\mathbf a,\mathbf b)=(\varphi\mathbf a,\varphi^{-1} \mathbf b) .\]
作用。于是外积 $\mathbf{ab}$ 是秩 $\leq 1$ 的 $n\times n$ 矩阵，并且在 $G$ 的作用下不变。因此，$(\mathbf a,\mathbf b)\mapsto \mathbf{ab}$ 定义了一个映射 $\mathfrak{M}_{\operatorname{B}}(T^*\C^n,G)\to M_{n\times n}(\C)$。矩映射 $\mu(\mathbf a,\mathbf b)=\mathbf{ba}$ 由点积给出，故 $\mu(\mathbf a,\mathbf b)=0$ 当且仅当 $\mathbf{ab}$ 幂零。于是，我们得到一个到秩 $\leq 1$ 的幂零矩阵空间 $\rkone$ 的映射 $\mathfrak{M}_{\operatorname{B}}(T^*\C^n,G)\to \rkone$。

过 $(\mathbf a,\mathbf b)$ 的 $G$-轨道是闭的，当且仅当 $\mathbf a$ 与 $\mathbf b$ 同时非零或同时为零；由此可以看出上述映射是双射。于是 $\mathfrak{M}_{\operatorname{B}}( T^*\C^n,G) \cong  \rkone$。
\end{ex}
\begin{ex}\label{ex:CZn}
设 $G\subset D$ 为行列式为 1 的对角矩阵构成的子群。我们仍可把 $T^*N$ 看作对 $(\mathbf a,\mathbf b)$。此时，矩映射条件保证对一切 $i,j$ 有 $a_ib_i=a_jb_j$，而闭轨道条件则保证：若对某个 $i$ 有 $a_i=0$，则对一切其他 $j$ 有 $a_j=0$；对 $b_*$ 同理。
通过乘以一个对角矩阵，我们可以进一步假设 $a_1=\cdots =a_n$ 且 $b_1=\cdots =b_n$。这定义了满射 $\C^2 \to \mathfrak{M}_{\operatorname{B}}$，把 $(x,y)$ 送到 $\mathbf{a}=(x,\dots, x)$ 与 $\mathbf{b}=(y,\dots,y)$。然而它不是单射：对角矩阵 $e^{2\pi ik/n} I$ 的行列式为 1，给出循环群 $\Z_n$ 在这个 $\C^2$ 的像上的作用。这与矩阵 $\operatorname{diag}(e^{2\pi ik/n},e^{-2\pi ik/n})$ 在 $\C^2$ 上的作用相吻合。于是 $\mathfrak{M}_{\operatorname{B}}(T^*\C^n,G)\cong \C^2/\mathbb{Z}_{n}$。
\end{ex}

其他环面群的超环面簇虽不那么常见、也更为复杂，却仍有很强的组合色彩；关于它们的问题通常可以归结为超平面配置的研究，正如环面簇可以用多面体来研究一样。值得注意的是，它们全都具有辛消解，可由 GIT 商、等价地由非零矩映射值处的超哈密顿约化（hyperhamiltonian reduction）构造。对上面的例 \ref{ex:TPn}，这个消解是 $T^*\mathbb{CP}^{n-1}$；对例 \ref{ex:CZn}，则是通过在奇点处反复爆破得到的唯一 crepant 消解。

\subsubsection{箭图规范理论}

这类约化最有名的例子是{\bf Nakajima 箭图簇}。当 $N$ 是箭图在固定向量空间上的表示空间时，它们就会出现。也就是说，我们固定一个有向图 $\Gamma$，以及一对向量 $\Bv,\Bw$，其分量由顶点集 $I$ 标号。群 $G=\prod_{i\in I} \operatorname{GL}(v_i;\C)$ 对每个 $i\in I$ 有表示 $\C^{v_i}$。我们要考虑的表示是
\[N= \Big(\bigoplus_{i\to j}\Hom(\C^{v_i},\C^{v_j})\Big)\oplus \Big(\bigoplus_{i\in I} \Hom(\C^{v_i},\C^{w_i})\Big).\]
我们希望得到的是左作用，因此 $(A,B)\in \operatorname{GL}(v_i;\C)\times \operatorname{GL}(v_j;\C)$ 作用在 $M\in \Hom(\C^{v_i},\C^{v_j})$ 上为 $BMA^{-1}$。我们称相应于这样选取的 $G$ 与 $N$ 的规范理论 $\mathcal{T}(T^*N,G)$ 为{\bf 箭图规范理论}，其希格斯分支 $\mathfrak{M}_{\operatorname{B}}(T^*N,G)$ 就是 Nakajima 箭图簇。它们是具有 Dynkin 图 $\Gamma$ 的 Kac--Moody 代数、最高权为 $\lambda$ 的表示的 $\mu$ 权空间的几何化身。它们由 $\alpha_i^{\vee}(\lambda)=w_i$（也就是说，$w_i$ 是它在与顶点 $i$ 对应的 $\mathfrak{sl}_2$ 下的权）以及 $\lambda-\mu=\sum v_i\alpha_i$（也就是说，$v_i$ 告诉我们从最高权空间往下走了多远）刻画。
\begin{yizhu}[Nakajima 箭图簇]
箭图簇是箭图表示空间上的超凯勒商：固定顶点维数 $\Bv$ 与框维数 $\Bw$，取 $T^*N$ 的矩映射零纤维，再模去 $G=\prod_{i\in I}\operatorname{GL}(v_i;\C)$ 的作用。这类簇由 Nakajima 引入，用来实现 Kac--Moody 代数最高权表示的权空间，故原文称之为「几何化身」；其中 $\Bw$ 记录最高权 $\lambda$ 在各单余根上的取值，$\Bv$ 记录 $\lambda$ 与权 $\mu$ 之差。
\end{yizhu}

我们可以把 $N$ 的点看作「带框箭图表示」，也就是如下配置：
\begin{enumerate}
    \item 给每个顶点 $i$ 一个向量空间 $\C^{v_i}$；
    \item 给每个顶点一个框映射（framing map）$\C^{v_i}\to \C^{w_i}$；
    \item 给每条有向边 $i\to j$ 一个线性映射 $\C^{v_i} \to \C^{v_j}$。
\end{enumerate}  
$G$ 的作用方式是改变 $\C^{v_i}$ 中的基，因此 $N$ 中的 $G$-轨道对应于带框箭图表示的同构类\footnote{使这一陈述成立所需的带框箭图表示范畴的定义，留作读者的练习。}。
物理学家通常会为每个顶点画两份：一份画成圆，里面填 $v_i$；一份画成方框，里面填 $w_i$。然后在第一份之间画出 $\Gamma$ 的边，再在同一个顶点的圆形副本与方形副本之间连边（标号为 0 的顶点不画）。
于是我们可以把 $T^*N$ 解释为双倍箭图（doubled quiver）的表示。

\begin{ex}\label{ex:Gr}
  若箭图只有一个顶点，于是只有一个 $v$ 与一个 $w$，则 $N=\Hom(\C^v,\C^w)$，而 $X=T^*N=\Hom(\C^v,\C^w)\oplus \Hom(\C^w,\C^v)$，也就是一对矩阵 $(A,B)$，其中 $A$ 为 $w\times v$、$B$ 为 $v\times w$。此时矩映射为
$\mu(A,B)=BA$。若 $v=1$，这就归结为例 \ref{ex:TPn}；更一般地，矩阵 $AB$ 在 $\operatorname{GL}(v)$ 的作用下不变，并给出同构
\[\mathfrak{M}_{\operatorname{B}}(T^*N,G)\cong \{C\in M_{w\times w}(\C) \mid C^2=0,\ \operatorname{rk}(C)\leq v\}.\]
\end{ex} 
 
其他重要的例子：
\begin{enumerate}
    \item[E1.] 这个箭图给出 $n\times n$ 幂零矩阵构成的空间：\begin{equation*}
\tikz{
	\node[draw, thick, circle, inner sep=6pt,fill=white] (a) at (-4.5,0) {$1$};
	\node[draw, thick, circle, inner sep=6pt,fill=white] (b) at (-3,0) {$2$};	
		\node[draw, thick, circle, inner sep=2pt,fill=white] (d) at (0,0) {$n-1$};
	\node[draw, thick,inner sep=12pt] (s) at (1.5,0){$n$};
	\node[inner sep=12pt,fill=white] (c) at (-1.5,0){$\cdots$};
		\draw[thick] (a) -- (b) ;
		\draw[thick] (b) -- (c) ;
		\draw[thick] (c) -- (d) ;
		\draw[thick] (d) -- (s) ;
}
\end{equation*}

\item[E2.] 这个箭图给出对称幂 $\Sym^n\C^2=(\C^2)^n/S_n$。
\begin{equation*}
\tikz{
		\node[draw, thick, circle, inner sep=6pt,fill=white] (d) at (0,0) {$n$};
	\node[draw, thick,inner sep=12pt] (s) at (2,0){$1$};
		\draw[thick] (d.north) to[out=90,in=0] (-.5,.8) to[out=180,in=90] (-1.25,0) to[out=-90,in=180] (-.5,-.8) to[out=0,in=-90] (d.south) ;
		\draw[thick] (d) -- (s) ;
}\end{equation*}

\end{enumerate}

这类构造还有许多变体，但对我们来说，它们在本文余下部分已足以充当主要例子。所有这些例子都具有辛消解，办法是把仿射商换成 GIT 商：秩 $\leq v$ 的矩阵空间由 $T^*\operatorname{Gr}(w,v)$ 消解（当 $v\leq w/2$ 时）；幂零锥由旗簇的余切丛消解（这是 Springer 消解的一个特例）；对称幂 $\Sym^n\C^2$ 则由 $\C^2$ 上 $n$ 个点的 Hilbert 概形消解。所有满足合理技术条件的 Nakajima 箭图簇都具有这样的消解。
% [原文 p10–p15]
